\chapter{Geometri dan Simetri}
\label{ch:gdl}

Kuliah geometric deep learning datang di semester terakhir, pada musim panas 2026. Materinya
mengikuti kursus Geometric Deep Learning dari Oxford yang dikembangkan Michael Bronstein dan
rekan-rekannya, dengan buku proto karya \citet{bronstein2021} sebagai acuan. Kuliah-kuliah sebelumnya
membahas arsitektur satu per satu. Kuliah ini menunjukkan bahwa konvolusi, DeepSets, jaringan graf,
dan transformer adalah anak-anak dari satu resep yang sama.

Resep itu bertumpu pada satu pengamatan. Di \cref{ch:data} kita melihat bahwa belajar di dimensi
tinggi hampir mustahil tanpa asumsi. Asumsi yang paling berharga ternyata adalah simetri: kita
sering tahu transformasi apa yang tidak mengubah jawaban. Kucing yang digeser tetap kucing.
Molekul yang diputar tetap molekul yang sama. Urutan penomoran atom dalam berkas tidak mengubah
sifat kimianya. Kalau pengetahuan seperti ini ditanamkan ke dalam arsitektur, ruang yang harus
dicari menyusut drastis.

\section{Sinyal, domain, dan grup}

Bahasa geometric deep learning dimulai dengan memisahkan dua hal: domain dan sinyal. Domain
$\Omega$ adalah tempat data hidup, misalnya grid piksel, himpunan titik, graf, atau permukaan.
Sinyal adalah fungsi $x:\Omega\to\R^c$ yang memberi setiap titik domain sebuah vektor fitur,
misalnya warna untuk setiap piksel atau sifat kimia untuk setiap atom. Ruang semua sinyal pada
sebuah domain, $\mathcal X(\Omega)$, adalah ruang vektor: dua sinyal bisa dijumlahkan dan dikalikan
skalar titik demi titik. Struktur linear inilah yang memungkinkan kita berbicara tentang lapisan
linear.

Simetri dirumuskan dengan \istilah{grup}. Grup adalah himpunan $G$ dengan operasi komposisi yang
memenuhi empat aksioma: hasil komposisi dua anggota tetap di dalam $G$, komposisinya asosiatif,
ada anggota identitas, dan setiap anggota punya invers. Rotasi bidang membentuk grup, begitu pula
pergeseran grid dan permutasi himpunan. Dua grup bisa memiliki struktur yang sama walau anggotanya
berbeda, dan kesamaan ini dinyatakan dengan \istilah{homomorphism}, pemetaan $\varphi$ yang menghormati
komposisi, $\varphi(gh)=\varphi(g)\varphi(h)$. Determinan adalah contohnya: ia memetakan matriks yang
bisa dibalik ke bilangan real tak nol, dan $\det(\bm A\bm B)=\det\bm A\det\bm B$. Fungsi eksponensial
memberi contoh yang lebih menarik, karena ia memetakan bilangan real dengan penjumlahan ke bilangan
positif dengan perkalian, $e^{a+b}=e^ae^b$, dan bisa dibalik dengan logaritma. Homomorfisme yang
bisa dibalik disebut isomorfisme, dan dua grup yang isomorfik secara struktural sama.

Grup bekerja pada domain melalui \istilah{group action}: setiap $g$ memindahkan titik $u$ ke $g.u$,
dengan identitas tidak memindahkan apa pun dan aksi dua anggota berturut-turut sama dengan aksi
komposisinya. Dari aksi pada domain, kita mendapat aksi pada sinyal,
\begin{equation}
\big(\rho(g)x\big)(u)=x\big(g^{-1}.u\big).
\label{eq:representasi}
\end{equation}
Pemetaan $g\mapsto\rho(g)$ ke operator linear pada ruang sinyal disebut \istilah{representasi}. Invers
pada \cref{eq:representasi} bukan hiasan. Untuk menggeser sebuah gambar ke kanan, nilai piksel di
posisi baru diambil dari posisi lama di sebelah kirinya. Tanpa invers, komposisi dua pergeseran
akan keluar dengan urutan terbalik, dan $\rho$ tidak lagi menjadi homomorfisme.

\section{Invariansi dan ekuivariansi}

Misalkan $\rho$ dan $\rho'$ representasi grup yang sama pada ruang input dan ruang output. Fungsi
$f$ disebut \istilah{invariant} kalau $f(\rho(g)x)=f(x)$ untuk setiap $g$, dan \istilah{equivariant} kalau
\begin{equation}
f\big(\rho(g)x\big)=\rho'(g)\,f(x)\qquad\text{untuk setiap }g\in G.
\end{equation}
Klasifikasi gambar seharusnya invarian terhadap pergeseran: labelnya tidak berubah. Segmentasi
seharusnya ekuivarian: kalau gambar digeser, peta segmentasinya ikut bergeser. Secara visual,
ekuivariansi berarti diagram berikut komut: menerapkan transformasi lalu $f$ memberi hasil yang sama
dengan menerapkan $f$ lalu transformasi padanannya. Sinyal-sinyal yang hanya berbeda karena
transformasi grup membentuk sebuah \istilah{orbit}. Peta ekuivarian memetakan setiap orbit di ruang
input ke sebuah orbit di ruang fitur, sehingga struktur simetrinya terbawa utuh.

Ada cara umum untuk membuat sembarang fungsi menjadi invarian terhadap grup yang terbatas, yaitu
merata-ratakannya atas seluruh grup,
\begin{equation}
(S_Gf)(x)=\frac1{|G|}\sum_{g\in G}f\big(\rho(g)x\big).
\end{equation}
Operator ini kadang disebut operator simetrisasi. Ambil himpunan dua elemen dan fungsi
$f(x_1,x_2)=x_1$ yang jelas tidak peduli pada simetri. Rata-rata atas dua permutasi memberi
$(x_1+x_2)/2$, yang invarian. Kalau fungsi target memang invarian, membatasi diri pada kelas fungsi
invarian tidak menambah galat aproksimasi, karena targetnya masih ada di dalam kelas. Galat
estimasi justru bisa turun, karena setiap contoh data kini ``mewakili'' seluruh orbitnya. Hanya
saja, kalau grupnya kecil dibanding dimensi masalah, galat estimasi tetap bisa terkena kutukan
dimensi. Simetri membantu, tetapi tidak membebaskan sepenuhnya.

\section{Simetri sebagai penghemat data}

Tugas pertama kuliah ini membuat kutukan dimensi dari \cref{ch:data} terasa dengan angka. Untuk dua titik acak di hiperkubus satuan
$[0,1]^n$, setiap koordinat menyumbang selisih kuadrat dengan rata-rata $1/6$, sehingga rata-rata kuadrat jaraknya $n/6$. Jarak tumbuh
seperti $\sqrt n$, sementara sebarannya relatif terhadap rata-rata justru menyempit, sehingga semua pasangan titik tampak sama jauhnya.

Tugas itu juga memperlihatkan apa yang dilakukan fitur yang tidak berguna. Dua kelas yang terpisah dengan baik dalam dua fitur,
misalnya dengan kuadrat jarak antarkelas sekitar 4 dan dalam kelas sekitar 1, menjadi sulit dipisahkan begitu ditambahi fitur derau.
Setiap fitur derau dengan simpangan baku $\sigma$ menambah rata-rata $2\sigma^2$ pada kuadrat jarak \emph{setiap} pasangan, sama untuk
pasangan sekelas maupun beda kelas. Seratus fitur derau dengan $\sigma=0{,}5$ menambah 50 pada keduanya, sehingga rasio jarak beda kelas
terhadap jarak sekelas turun dari 4 menjadi $54/51\approx1{,}06$. Tetangga terdekat praktis menjadi acak. Inilah sebabnya
$k$-nearest neighbour dan pengelompokan peka sekali terhadap pemilihan fitur (\cref{ch:unsup}).

Bagian kedua tugas memakai integrasi Monte Carlo (\cref{ch:fiskom}) untuk fungsi berupa gumpalan Gauss yang tersusun simetris di
dalam hiperkubus, dan meminta kami memanfaatkan simetrinya. Kalau fungsi tidak berubah ketika setiap koordinat dicerminkan dan ketika
koordinat-koordinat ditukar urutannya, integralnya atas seluruh kubus sama dengan integral atas satu ``potongan dasar'' dikali jumlah
transformasi itu. Dalam $d$ dimensi ada $2^d$ pencerminan dan $d!$ permutasi. Untuk $d=8$, itu $2^8\cdot8!\approx10^7$ potongan yang
identik. Sampel cukup ditarik dari satu potongan, yang volumenya sepersepuluh juta volume kubus, tanpa kehilangan informasi apa pun.

Pelajaran yang sama berlaku untuk belajar, bukan hanya untuk integrasi. Model yang invarian terhadap sebuah grup tidak perlu melihat
setiap versi yang ditransformasi dari setiap contoh, karena satu contoh sudah mewakili seluruh orbitnya. Augmentasi data mencoba
mencapai hal yang sama dengan memperlihatkan versi-versi itu secara eksplisit. Arsitektur yang ekuivarian mencapainya dengan
konstruksi, dan untuk grup yang besar, selisihnya bisa berorde jutaan, seperti contoh di atas.

\section{Skala dan cetak biru}

Simetri saja belum cukup. Prior kedua adalah \istilah{scale separation}: banyak sinyal bisa
dipahami secara bertingkat. Detail lokal memengaruhi detail lokal di sekitarnya, dan struktur
besar bisa diringkas dari versi sinyal yang lebih kasar. Dua bahan dasarnya adalah operator lokal
yang bekerja di lingkungan kecil, dan operasi \istilah{coarse graining} yang meringkas sinyal ke
domain yang lebih kasar, seperti pooling pada gambar. Gabungan simetri dan skala lebih kuat dari
masing-masing, karena fungsi yang bekerja secara lokal dan bertingkat hanya perlu belajar
interaksi jarak pendek di setiap tingkat.

Dari kedua prior ini, \citet{bronstein2021} merumuskan sebuah cetak biru. Jaringan disusun dari
lapisan linear yang ekuivarian terhadap $G$ (sebut $B$), nonlinearitas yang diterapkan titik demi
titik ($\sigma$), pooling lokal yang mengasarkan domain ($P$), dan di ujungnya pooling global yang
invarian ($A$), sehingga
\begin{equation}
f=A\circ\sigma\circ B_J\circ P_{J-1}\circ\cdots\circ P_1\circ\sigma\circ B_1.
\end{equation}
CNN adalah cetak biru ini pada grid dengan grup translasi. DeepSets adalah cetak biru pada
himpunan dengan grup permutasi. Jaringan graf adalah cetak biru pada graf. Kita akan melihat
ketiganya satu per satu.

Ada satu kehalusan lagi. Gambar nyata jarang bertransformasi persis menurut grup. Wajah yang
sedikit tersenyum bukan rotasi atau pergeseran dari wajah yang datar, melainkan deformasi kecil.
Representasi yang baik sebaiknya \emph{hampir} invarian terhadap transformasi yang \emph{hampir}
berada di dalam grup: perubahan keluarannya sebanding dengan seberapa jauh deformasi itu dari
transformasi grup. Sifat ini disebut stabilitas terhadap deformasi, dan ia menjadi ukuran kualitas
representasi yang akan muncul lagi pada wavelet dan pada filter spektral.

\section{Himpunan}

Domain paling sederhana adalah himpunan tanpa struktur. Simetrinya adalah grup permutasi: penomoran
elemen tidak boleh berpengaruh. Sinyal pada himpunan $n$ elemen bisa ditulis sebagai matriks fitur
$\mX\in\R^{n\times c}$, satu baris per elemen. Kalau elemen dipermutasi, baris-baris matriks ikut
dipermutasi, $\mX\mapsto\bm P\mX$, dengan $\bm P$ matriks permutasi yang berisi tepat satu angka
satu di setiap baris dan kolom.

Peta linear yang ekuivarian terhadap semua permutasi ternyata hanya punya dua derajat kebebasan
per pasangan kanal: satu bagian yang bekerja pada setiap elemen sendiri, dan satu bagian yang
melihat jumlah seluruh elemen,
\begin{equation}
F(\mX)=\mX\bm\Lambda+\tfrac1n\bm 1\bm 1^\top\mX\,\bm\Gamma.
\end{equation}
DeepSets \citep{zaheer2017} membangun fungsi invarian dari bagian-bagian ini:
$f(\mX)=\phi\big(\sum_i\psi(\vx_i)\big)$. Fungsi $\psi$ dipakai bersama untuk setiap elemen dan
bersifat ekuivarian, penjumlahan adalah pooling yang invarian, dan $\phi$ bekerja pada ringkasan
global. Kedua fungsi dipelajari, sedangkan penjumlahannya tetap.

\section{Graf}

Graf menambahkan satu bahan ke himpunan: hubungan antar elemen. Graf tak berarah dengan atribut
simpul bisa ditulis sebagai matriks fitur $\mX$ dan matriks ketetanggaan $\bm A$, dengan
$a_{uv}=1$ kalau simpul $u$ dan $v$ terhubung. Grup simetrinya tetap permutasi, tetapi sekarang
permutasi bekerja pada keduanya: $\mX\mapsto\bm P\mX$ dan $\bm A\mapsto\bm P\bm A\bm P^\top$. Fungsi
tingkat graf harus invarian, $f(\bm P\mX,\bm P\bm A\bm P^\top)=f(\mX,\bm A)$, dan fungsi tingkat
simpul harus ekuivarian, $F(\bm P\mX,\bm P\bm A\bm P^\top)=\bm P\,F(\mX,\bm A)$.

\subsection{Message passing}

Graph neural network modern hampir semuanya mengikuti pola \istilah{message passing}
\citep{gilmer2017,battaglia2018}. Setiap simpul $u$ mengumpulkan pesan dari tetangganya, lalu
memperbarui keadaannya:
\begin{equation}
\vh_u=\phi\Big(\vx_u,\ \bigoplus_{v\in\mathcal N_u}\psi(\vx_u,\vx_v)\Big),
\end{equation}
dengan $\bigoplus$ operator agregasi yang tidak peduli urutan, seperti jumlah, rata-rata, atau
maksimum. Bedanya dengan DeepSets ada di sini: setiap simpul tidak lagi meringkas seluruh himpunan,
tetapi hanya tetangganya, sehingga sisi graf menentukan siapa berbicara dengan siapa.

Pembagian yang dipakai di kuliah membedakan tiga rasa. Pada varian konvolusional, pesan dari
tetangga dikalikan bobot tetap yang bergantung pada struktur graf, seperti pada GCN
\citep{kipf2017}. Pada varian attentional, bobotnya dihitung dari fitur kedua simpul, seperti pada
GAT \citep{velickovic2018}. Pada varian message passing yang paling umum, pesan $\psi$ adalah fungsi
sembarang dari kedua simpul, dan kedua varian sebelumnya adalah kasus khususnya.

Sebagai contoh kecil, ambil graf garis dengan tiga simpul, $1-2-3$, dengan fitur skalar
$x=(1,2,3)$. Dengan aturan $h_u=x_u+\sum_{v\in\mathcal N_u}x_v$, simpul pertama mendapat $1+2=3$,
simpul kedua $2+1+3=6$, dan simpul ketiga $3+2=5$. Kalau penomoran dibalik menjadi $3-2-1$, fitur
menjadi $(3,2,1)$ dan hasilnya $(5,6,3)$, yaitu hasil semula dalam urutan baru. Itulah
ekuivariansi.

Contoh ini bukan kebetulan. Dalam bentuk matriks, aturan tadi adalah $\vh=(\bm I+\bm A)\vx$. Kalau simpul
dipermutasi dengan $\bm P$, input menjadi $\bm P\vx$ dan ketetanggaan menjadi $\bm P\bm A\bm P^\top$, sehingga
keluarannya $(\bm I+\bm P\bm A\bm P^\top)\bm P\vx=\bm P\vx+\bm P\bm A\vx=\bm P(\bm I+\bm A)\vx=\bm P\vh$,
karena $\bm P^\top\bm P=\bm I$. Argumen yang sama berlaku untuk message passing umum: fungsi $\psi$ dan
$\phi$ yang dipakai bersama di semua simpul tidak peduli nomor simpulnya, dan agregasi atas tetangga
tidak peduli urutan tetangganya. Satu-satunya hal yang bisa membuat jaringan bergantung pada penomoran
adalah memberi setiap simpul bobot pribadinya sendiri, dan itulah yang dihindari.

Teman yang mengabarkan berita di sebuah desa adalah gambaran yang cocok. Setiap pagi, setiap
orang mendengar kabar dari tetangga sebelah rumahnya dan memperbarui pengetahuannya. Setelah satu
pagi, kabar dari rumah sebelah sudah sampai. Setelah tiga pagi, kabar dari rumah yang berjarak
tiga sudah sampai. Satu lapisan message passing adalah satu pagi.

\subsection{Seberapa kuat jaringan graf?}

Bisakah jaringan graf membedakan dua graf yang tidak isomorfik? Pertanyaan ini terhubung dengan
uji klasik \istilah{Weisfeiler--Leman} (WL). Uji ini mewarnai setiap simpul, lalu berulang kali
memberi warna baru berdasarkan warna lama simpul itu dan multihimpunan warna tetangganya. Kalau
histogram warna dua graf berbeda, graf itu pasti tidak isomorfik. Kalau sama, kita tidak bisa
menyimpulkan apa-apa. \citet{xu2019} dan \citet{morris2019} menunjukkan bahwa message passing paling
banter sekuat WL, dan mencapai kekuatan itu kalau agregasinya injektif terhadap multihimpunan.
Penjumlahan memenuhi syarat ini, sedangkan rata-rata dan maksimum tidak, karena keduanya tidak bisa
membedakan, misalnya, satu tetangga berfitur satu dari dua tetangga berfitur satu.

Batas WL mudah diperlihatkan. Bandingkan dua segitiga yang terpisah dengan satu lingkaran enam
simpul. Keduanya punya enam simpul, dan setiap simpul punya tepat dua tetangga. Pada setiap iterasi
WL, semua simpul menerima informasi yang sama, sehingga semua warna tetap identik dan uji tidak
bisa membedakan keduanya, padahal yang satu berisi segitiga dan yang lain tidak. Jaringan message
passing biasa pun gagal di sini.

Ada beberapa cara melampaui batas ini. \istilah{Positional encoding} memberi setiap simpul ``alamat'',
misalnya dari vektor eigen Laplacian graf atau dari langkah acak, sehingga simpul yang
strukturnya sama tetap bisa dibedakan. Jaringan substruktur menghitung kemunculan pola kecil,
seperti segitiga atau siklus, di sekitar setiap simpul. Subgraph GNN menjalankan jaringan pada
banyak versi graf yang sedikit diubah, misalnya dengan satu simpul dihapus, lalu menggabungkan
hasilnya. Gagasan ini berkerabat dengan \istilah{reconstruction conjecture}, dugaan lama dalam teori
graf bahwa sebuah graf dengan setidaknya tiga simpul ditentukan sepenuhnya oleh kumpulan subgraf
yang masing-masing kehilangan satu simpul.

Transformer pun bisa dilihat sebagai message passing pada graf lengkap, di mana setiap token adalah
simpul yang menerima pesan dari semua simpul lain dengan bobot attention. Tanpa positional
encoding, attention memang ekuivarian terhadap permutasi. Dengan positional encoding seperti yang
dipakai dalam praktik, sifat itu sengaja dilepas, karena urutan kata memang bermakna.

\subsection{Oversquashing}

Untuk menangkap interaksi jarak jauh, kita menambah lapisan. Tetapi jumlah simpul dalam radius $r$
bisa tumbuh secara eksponensial, sementara semua informasinya harus dimampatkan ke vektor
berukuran tetap. Gejala ini disebut \istilah{oversquashing} \citep{alon2021}. Secara formal, kepekaan
keadaan simpul $u$ setelah $r$ lapisan terhadap fitur awal simpul $v$,
$\|\partial\vh_u^{(r)}/\partial\vx_v\|$, dibatasi oleh entri matriks ketetanggaan ternormalisasi
yang dipangkatkan, $(\hat{\bm A}^r)_{uv}$, dikalikan konstanta yang bergantung pada bobot dan
nonlinearitas. Lebar jaringan dan kedalaman ikut berperan, tetapi faktor yang paling menentukan
adalah topologi graf: leher botol, yaitu sedikit sisi yang menghubungkan dua kelompok besar, dan
daerah dengan kelengkungan negatif \citep{topping2022}. Kabar dari desa sebelah harus lewat satu
jembatan sempit, dan di jembatan itu sebagian besar berita hilang.

\section{Grid}

Grid adalah kasus yang paling kita kenal. Grid satu dimensi bisa dilihat sebagai graf berbentuk
lingkaran, kalau tepinya disambung secara periodik. Grup simetrinya adalah pergeseran siklik.
Gambar $\bm I$ yang digeser satu piksel ke kanan bisa ditulis $\bm I'=\bm A\bm I\bm B$, dengan
$\bm A$ matriks identitas karena baris tidak berubah, dan $\bm B$ matriks pergeseran siklik yang
memindahkan setiap kolom satu tempat. Pergeseran ke bawah memakai matriks pergeseran di sisi kiri.

Peta linear yang ekuivarian terhadap pergeseran adalah tepat peta yang komut dengan operator
geser $\bm S$, yaitu $\bm C\bm S=\bm S\bm C$. Matriks seperti ini adalah \istilah{circulant matrix}:
setiap barisnya adalah baris sebelumnya yang digeser satu tempat. Mengalikan dengan matriks
sirkulan sama dengan konvolusi. Inilah alasan mendalam mengapa konvolusi muncul di CNN: ia
bukan sekadar pilihan desain, melainkan satu-satunya operasi linear yang menghormati simetri geser.
Semua matriks sirkulan punya vektor eigen yang sama, yaitu basis Fourier, sehingga transformasi
Fourier diskret mendiagonalkan semuanya sekaligus. Akibatnya, konvolusi di ruang menjadi perkalian
titik demi titik di frekuensi, dan dengan FFT \citep{cooley1965}, konvolusi panjang bisa dihitung
dalam $\mathcal O(n\log n)$ operasi.

Fourier punya keterbatasan. Basisnya tersebar di seluruh ruang, sehingga lokalisasi sempurna di
frekuensi berarti tidak ada lokalisasi di ruang. Prinsip ketidakpastian membatasi seberapa sempit
sebuah fungsi bisa sekaligus di kedua domain. Transformasi Fourier berjendela memotong sinyal
menjadi potongan pendek, tetapi lebar jendelanya sama untuk semua frekuensi. Lebih penting lagi,
modulus Fourier tidak stabil terhadap dilatasi kecil. Meregangkan sinyal sedikit menggeser
frekuensi rendah sedikit, tetapi menggeser frekuensi tinggi jauh, sehingga dua sinyal yang hampir
sama bisa punya spektrum frekuensi tinggi yang berbeda.

Wavelet menyelesaikan ini dengan memakai satu fungsi induk yang digeser dan diskalakan. Di
frekuensi tinggi, wavelet sempit di ruang dan lebar di frekuensi. Di frekuensi rendah, sebaliknya.
Wavelet Haar, misalnya, adalah tangga naik-turun yang cocok sekali untuk mendeteksi tepi dan
struktur lokal. \istilah{Scattering transform} \citep{mallat2012} menumpuk tiga bahan: konvolusi
wavelet, modulus, dan rata-rata. Modulus adalah nonlinearitas yang krusial. Tanpanya, rata-rata
dari koefisien wavelet akan bernilai nol, karena wavelet berosilasi. Modulus memindahkan energi
frekuensi tinggi ke frekuensi rendah, sehingga informasinya selamat dari rata-rata. Hasilnya adalah
representasi yang invarian terhadap pergeseran dan stabil terhadap deformasi, sebuah CNN yang
filternya tidak dipelajari tetapi dirancang, dan dalam banyak hal menjadi model teoretis untuk
memahami CNN.

\section{Konvolusi grup}

Konvolusi bisa dipandang sebagai ``transformasikan lalu cocokkan'': geser filter ke setiap posisi,
lalu ukur kecocokannya dengan sinyal. Kalau grupnya lebih besar dari pergeseran, misalnya
pergeseran ditambah rotasi, kita bisa melakukan hal yang sama untuk setiap anggota grup
\citep{cohen2016}:
\begin{equation}
(x\star\psi)(g)=\int_\Omega x(u)\,\psi(g^{-1}.u)\dd u.
\end{equation}
Keluarannya kini adalah fungsi pada grup, bukan pada domain asal. Pada grup p4, yang terdiri atas
pergeseran dan empat rotasi kelipatan sembilan puluh derajat, lapisan pertama mengubah gambar
menjadi empat peta fitur, satu untuk setiap orientasi filter. Lapisan berikutnya bekerja pada
fungsi di grup itu sendiri, sehingga filternya juga punya dimensi orientasi. Kalau gambar input
diputar sembilan puluh derajat, setiap peta fitur ikut berputar, dan keempat orientasi bergeser
secara siklis.

Konsep pendukungnya adalah \istilah{homogeneous space}. Grup bekerja secara transitif pada sebuah
domain kalau setiap titik bisa dicapai dari titik lain oleh suatu anggota grup. Stabilizer sebuah
titik adalah subgrup yang membiarkan titik itu di tempatnya. Bola $S^2$ di bawah rotasi tiga
dimensi adalah contoh homogeneous space, dan stabilizer kutub utaranya adalah rotasi di sekitar
sumbu yang melewati kutub. Bidang di bawah gerak kaku, yaitu pergeseran dan rotasi, adalah contoh
lainnya, dengan stabilizer setiap titik berupa rotasi di sekitar titik itu.

\section{Manifold dan mesh}

Banyak data hidup di permukaan melengkung: tubuh manusia, sayap pesawat, permukaan bumi. Manifold
adalah model matematisnya. Secara informal, manifold adalah ruang yang di sekitar setiap titik
tampak seperti ruang Euklides. \istilah{Chart} adalah peta lokal yang meratakan sepotong manifold,
seperti halaman atlas yang meratakan sepotong bumi. Di setiap titik ada \istilah{tangent space},
ruang semua arah yang mungkin di titik itu. Untuk mengukur panjang dan sudut vektor singgung,
kita butuh \istilah{Riemannian metric}, sebuah hasil kali dalam yang bervariasi mulus dari titik ke
titik. Dengan metrik, kita bisa mendefinisikan panjang kurva, \istilah{geodesic} sebagai jalur
terpendek, dan jarak geodesik.

Besaran yang hanya bergantung pada metrik disebut intrinsik. Ia tidak berubah kalau permukaan
ditekuk tanpa diregangkan. Selembar kertas yang digulung menjadi silinder punya geometri intrinsik
yang sama dengan kertas datar, sehingga jarak geodesik antara dua titik di atasnya tidak berubah.
Deformasi yang mempertahankan metrik disebut isometri. Berbeda dari homogeneous space, manifold
umumnya tidak punya grup simetri global, sehingga stabilitas terhadap deformasi menjadi lebih
penting daripada ekuivariansi yang persis.

Kalkulus pada manifold bertumpu pada tiga operator. Gradien $\nabla f$ sebuah medan skalar menunjuk
ke arah kenaikan tercepat dan didefinisikan melalui metrik, sehingga intrinsik. Divergensi sebuah
medan vektor mengukur sumber dan sumur lokal. Keduanya saling adjoin, $\langle\nabla f,X\rangle=-\langle f,\operatorname{div}X\rangle$,
persis seperti integrasi parsial. Gabungannya adalah \istilah{Laplace--Beltrami operator},
$\Delta f=-\operatorname{div}\nabla f$, yang mengukur seberapa jauh nilai fungsi di sebuah titik
berbeda dari rata-rata di sekitarnya. Laplacian ini simetris dan semidefinit positif, sehingga punya
basis fungsi eigen yang ortogonal, generalisasi alami dari basis Fourier.

Dengan basis itu, konvolusi bisa didefinisikan secara spektral: transformasikan sinyal ke basis
eigen Laplacian, kalikan dengan filter di setiap frekuensi, lalu kembalikan \citep{bruna2014}.
Konvolusi spektral ini intrinsik, tetapi naif. Basis eigen bisa berubah drastis karena deformasi
kecil, terutama untuk frekuensi tinggi, sehingga filter yang dipelajari pada satu bentuk tidak
berpindah ke bentuk lain. Jalan keluarnya adalah filter yang mulus terhadap frekuensi, misalnya
polinomial dalam Laplacian. Filter polinomial berderajat $K$ hanya melihat tetangga dalam radius
$K$, sehingga lokal dan stabil, dan ChebNet \citep{defferrard2016} memakai polinomial Chebyshev untuk
tujuan ini.

Dalam praktik, manifold didiskretisasi menjadi mesh segitiga atau graf. Laplacian graf yang paling
sederhana adalah $\bm L=\bm D-\bm A$, dengan $\bm D$ matriks diagonal derajat simpul, sehingga
\begin{equation}
(\bm L\vx)_u=\sum_{v\in\mathcal N_u}a_{uv}\,(x_u-x_v).
\end{equation}
Pada mesh, bobot $a_{uv}$ dipilih dari geometri segitiga, misalnya dengan rumus kotangen, agar
Laplacian diskret mendekati Laplace--Beltrami yang kontinu. Bentuk di atas sudah memperlihatkan
bahwa Laplacian adalah sebuah langkah message passing: setiap simpul menjumlahkan selisih dirinya
dengan tetangganya. Konvolusi spektral, jaringan graf, dan operator diferensial pada permukaan
ternyata adalah satu keluarga.

\section{Simetri di luar kelas}

Cetak biru ini muncul di banyak bidang dengan nama yang berbeda. Molekul adalah graf atom, dan sifatnya
tidak berubah kalau molekul diputar atau atom-atomnya dinomori ulang, sehingga jaringan graf dan jaringan
ekuivarian menjadi standar dalam kimia komputasi (\cref{ch:hayati}). AlphaFold menempatkan residu protein
sebagai benda kaku di ruang dengan attention yang tidak peduli rotasi global (\cref{ch:protein}). Awan titik
lidar pada kendaraan otonom adalah himpunan tanpa urutan, dan PointNet adalah DeepSets untuk awan titik
(\cref{ch:robot}). Pengguna dan item dalam sistem rekomendasi membentuk graf bipartit, dan message passing
di atasnya menyebarkan selera (\cref{ch:rekomendasi}).

Bagi simulasi, bab ini barangkali yang paling penting di Bagian II. Hukum fisika punya simetri:
persamaan Navier--Stokes tidak berubah kalau sistem digeser atau diputar, dan gaya antar partikel
tidak bergantung pada urutan penomoran partikel. Model pengganti yang tidak menghormati simetri
ini harus mempelajarinya dari data, dan sering gagal di luar data latihnya. Jaringan graf yang
ekuivarian terhadap rotasi dan translasi tiga dimensi, seperti EGNN \citep{satorras2021},
membangunnya langsung ke dalam arsitektur.

Mesh simulasi pada dasarnya adalah graf. Metode volume hingga, yang akan kita bahas di
\cref{ch:numerik}, menghitung perubahan di setiap sel dari fluks yang melewati sisi-sisinya ke sel
tetangga. Bentuknya persis message passing: setiap sel mengumpulkan ``pesan'' fluks dari
tetangganya, lalu memperbarui keadaannya. Tidak mengherankan bahwa simulator yang dipelajari
dengan jaringan graf menjadi salah satu keberhasilan terbesar di bidang ini. Graph Network-based
Simulator \citep{sanchez2020} meniru dinamika partikel cairan dan pasir, MeshGraphNets
\citep{pfaff2021} bekerja pada mesh simulasi aerodinamika dan struktur, dan Message Passing Neural
PDE Solvers \citep{brandstetter2022} memakai message passing sebagai solver PDE umum. GraphCast
\citep{lam2023} membawa gagasan yang sama ke prakiraan cuaca global, dengan graf di atas mesh
ikosahedral yang membungkus bumi. Semuanya dibahas lebih dalam di \cref{ch:operator}.

Laplacian graf juga punya makna fisik yang langsung. Persamaan difusi panas pada graf,
$\dot{\vx}=-\kappa\bm L\vx$, didiskretisasi dengan Euler eksplisit menjadi
$\vx^{n+1}=\vx^n-\kappa\Delta t\,\bm L\vx^n$. Satu langkah waktu difusi adalah satu lapisan message
passing dengan bobot yang sudah ditentukan. Pandangan ini juga menjelaskan gejala
\istilah{oversmoothing} pada jaringan graf yang terlalu dalam: seperti panas yang lama-lama merata,
fitur semua simpul menjadi semakin mirip.

\section{Perkembangan riset}

Di dunia atom dan molekul, jaringan ekuivarian kini menjadi alat standar. MACE \citep{batatia2022} memakai pesan dengan
orde badan yang lebih tinggi: alih-alih hanya pasangan atom, setiap pesan sudah memuat interaksi sampai empat atom
sekaligus. Dengan begitu, dua langkah message passing sudah cukup, sehingga model lebih cepat dan lebih mudah
diparalelkan, dengan akurasi medan gaya yang setara atau lebih baik dari model sebelumnya. Equiformer
\citep{liao2023equiformer} membawa transformer ke graf atom tiga dimensi dengan mengganti operasi transformer dengan versi
ekuivarian berbasis representasi tak tereduksi dan hasil kali tensor, sesuai teori grup di awal bab ini.

Untuk graf umum, GraphGPS \citep{rampasek2022} memisahkan dua mekanisme: message passing lokal di sepanjang sisi graf,
dan attention global antar semua simpul, ditambah encoding posisi dan struktur. Pemisahan ini membuat biayanya linear
terhadap jumlah simpul dan sisi, sekaligus menjadi salah satu jawaban atas masalah oversquashing, karena informasi bisa
melompat lewat attention global. Geometric Algebra Transformer \citep{brehmer2023} mewakili titik, arah, translasi, dan
rotasi sebagai elemen aljabar Clifford proyektif berdimensi 16, sehingga satu arsitektur ekuivarian terhadap E(3) bisa
menangani banyak jenis objek geometris. Contoh penerapannya mencakup simulasi banyak benda, perencanaan gerak robot,
dan estimasi tegangan geser dinding pada mesh arteri, sebuah masalah aliran darah yang menyambung langsung ke
\cref{ch:operator}.

\sumberbab{Bab ini mengikuti kuliah Deep Learning: Geometric Techniques, yang memakai materi kursus
Geometric Deep Learning dari Oxford dan daftar topik ujiannya: kutukan dimensi, prior geometris,
himpunan, graf, grid, konvolusi grup, manifold, dan konvolusi spektral. Acuan utama adalah
\citet{bronstein2021}, dengan latar belakang di \citet{bronstein2017}.}
